\noindent\textit{Source and counted unit.} Yvain Bruned and Vladimir Dotsenko, \emph{Novikov algebras and multi-indices in regularity structures}, Annales Henri Lebesgue 9 (2026), 791--814, \url{https://doi.org/10.5802/ahl.275}. Published under CC BY 4.0: \url{https://creativecommons.org/licenses/by/4.0/}.
\par\noindent The counted unit is the whole Theorem 3.3 (free multi-Novikov realization). The canonical-structure proof, populated-monomial and dimension bounds, ordered and normal reductions, and both native rotation-correspondence diagrams are support for that one unit. Sections 2 and 3 below preserve the original numbering; no helper proposition is separately counted.
\par\noindent\textit{Editorial boundary.} The original author mathematics is retained, including its notation and wording. The journal class and branding are replaced by a standard wrapper. The seven cited bibliography entries are transcribed from the published bibliography, pp.812--813, with original printed labels; no missing BibTeX file is guessed. Knuth's rotation correspondence and elementary free differential-algebra/graded dimension facts are standard prerequisites; the current multi-operation reductions remain fully reproduced.
\par\noindent The source's finite-set working assumption and its infinite-set grading footnote are retained. In particular the source uses ``right-leaning'' within the proof of Proposition 3.8 after defining ``left-leaning'', and its Proposition 3.4 proof uses the symbols $(f_0,x_0)$ and $(x,m_1)$. These are retained as printed, not silently repaired or independently certified.
